% ══════════════════════════════════════════════════════════════════════════════
\section{Introduction}
\label{sec:intro}
% ══════════════════════════════════════════════════════════════════════════════

Between 2010 and 2024, the Sahel belt spanning Burkina Faso, Mali, and Niger experienced a more than tenfold increase in ACLED-recorded violent incidents, with jihadist groups, inter-communal militias, and state forces contesting control of productive agricultural land during and after the harvest season. Yet in the same years, satellite imagery recorded steadily improving vegetation conditions across the region, suggesting that violence and agricultural abundance were not moving in opposite directions. This pattern---better harvests coinciding with more conflict---is puzzling under the standard assumption that agricultural productivity pacifies rural populations by raising the opportunity cost of violence. It motivates the central question of this paper: does predicted crop yield reduce conflict, and does the answer depend on whether armed groups can appropriate the harvest?

The empirical literature on income shocks and conflict has converged on rainfall as a workhorse instrument for agricultural productivity, exploiting year-to-year variation in precipitation to estimate the effect of income on violence \citep{MiguelSatyanathSergenti2004}. This strategy is limited in two important respects. First, rainfall affects conflict through multiple channels beyond crop income---including heat stress, displacement, and pastoral resource competition---making it difficult to isolate the agricultural productivity mechanism specifically. Second, even when the goal is to measure an agricultural shock, rainfall maps imperfectly onto yields: the same precipitation deviation produces very different output depending on soil quality, crop variety, irrigation access, and management practices. Studies relying on climate proxies therefore estimate a mixture of effects rather than the income channel of interest. A direct measure of predicted crop productivity, constructed from the full set of agronomic inputs, offers a sharper test of the theory.

This paper constructs such a measure. I train an XGBoost model on 3,363 GROW-Africa administrative maize yield records from 204 Admin-2 units spanning 2000--2023 and 210 satellite and environmental features---vegetation indices (EVI, NDVI, GCVI), Sentinel-1 SAR backscatter, CHIRPS monthly precipitation, ERA5 growing and killing degree days, iSDA soil properties, and SRTM terrain covariates---to predict district-level maize yields for 1,338 GADM Admin-2 districts across seven Sub-Saharan African countries each year from 2010 to 2024. The model achieves a spatial leave-one-admin1-out $R^2$ of 0.35. Because these predictions are constructed entirely from exogenous environmental inputs, they are not mechanically contaminated by local conflict outcomes, addressing the primary endogeneity concern with observed yields. I merge the predictions with geocoded ACLED conflict event data and estimate two-way fixed-effects regressions on an 18,096 observation Admin-2~$\times$~year panel, controlling for district population with Admin-2 and year fixed effects and standard errors clustered at the Admin-2 level.

The main finding is a negative and precisely estimated yield--conflict elasticity. In the preferred TWFE OLS log-log specification, a 10\% increase in ML-predicted maize yield reduces log conflict events in the subsequent 3-month post-harvest window by 6.7\% ($\hat{\beta} = -0.666$, $\text{SE} = 0.153$, $p < 0.001$). This effect survives a spatial lag control for neighbour conflict ($\hat{\beta} = -0.377$, $p < 0.001$), strengthens at longer horizons (12-month: $\hat{\beta} = -1.290$, $p < 0.001$), and is concentrated in violence against civilians and battles rather than in riots, which is inconsistent with a broad economic grievance channel and consistent with armed-actor responses to changes in the returns to violence.

The paper's most novel contribution is the discovery that this aggregate pacifying effect conceals a fundamental reversal across agro-ecological zones. In Non-Sahel countries---Ethiopia, Malawi, Nigeria, and Tanzania---the yield elasticity is negative and significant ($\hat{\beta} = -0.467$, $p < 0.001$): higher predicted yields reduce conflict, consistent with the opportunity-cost mechanism. In Sahel countries---Burkina Faso, Mali, and Niger---where jihadist and inter-communal armed groups systematically extract agricultural surplus through coercive market taxation, cattle raiding, and territorial control of productive zones, the relationship reverses: the full-sample interaction model yields a Sahel differential of $+3.513$ ($p < 0.001$) and a net Sahel elasticity of $+2.383$, indicating that better harvests attract predation and increase conflict. This pattern is consistent with the resource-appropriation mechanism formalised by \citet{Grossman1991} and \citet{DalBoDalBo2011}, and mirrors evidence from \citet{DubeVargas2013} that positive shocks to easily appropriable commodities increase conflict. The sign reversal between zones is the paper's central empirical result.

This paper makes three contributions. First, relative to the rainfall-instrument literature, it provides a more direct test of the agricultural income channel by measuring predicted crop productivity rather than the upstream climatic inputs that produce it, using a validated machine-learning model rather than a meteorological proxy. Second, relative to subnational Africa studies such as \citet{HarariLaFerrara2018}, it exploits a richer multi-country panel and moves beyond drought indices to a yield-specific measure, allowing separation of the agronomic channel from other climate pathways. Third, it provides the first large-scale subnational evidence that the sign of the yield--conflict relationship is reversed in agro-pastoral Sahel environments relative to higher-rainfall Non-Sahel ones, demonstrating that aggregate estimates are not informative about the direction of the effect in any specific context.

The remainder of the paper proceeds as follows. Section~\ref{sec:theory} develops the theoretical framework. Section~\ref{sec:litreview} reviews the empirical literature. Section~\ref{sec:data} describes the data and panel construction. Section~\ref{sec:empirical} presents the empirical strategy. Section~\ref{sec:results} reports the main results and robustness checks. Section~\ref{sec:discussion} discusses mechanisms and limitations.
